%\newpage
\section{Widerstands-Messung}
Jeder Widerstand wird mit vier verschiedenen Messmethoden in einer Stromrichtung vermessen.
Der gleiche Widerstand wird auch mit den gleichen vier Messmethoden in die andere Stromrichtung vermessen.
Die Messung in die Gegenrichtung wird nur f\"ur die Erkennung auf Widerstand benutzt.
Wenn die Abweichung dieser beiden Messungen zu gro"s ist, ist es kein Widerstand.

\subsection{Widerstandsmessung mit den 680-Ohm-Widerst\"anden}
Die Messung des unbekannten Widerstandes Rx wird in zwei verschiedenen Wegen mit den \(680\Omega\)-Pr\"azisionswiderst\"anden
 durchgef\"uhrt.
Das Schaltbild dieser Messungen mit Testpin 1 (TP1) und Testpin 3 (TP3) werden vereinfacht in Abbildung~\ref{fig:RL1mes} und Abbildung~\ref{fig:RL2mes} als ein Beispiel von den sechs Kombinationsm\"oglichkeiten gezeigt.

\begin{figure}[H]
\centering
\includegraphics[]{../FIG/ResistormessL1.eps}
\caption{Messung Type 1 mit \(680\Omega\) }
\label{fig:RL1mes}
\end{figure}

\begin{figure}[H]
 \centering
 \includegraphics[]{../FIG/ResistormessL2.eps}
 \caption{Messung Type 2 mit \(680\Omega\) }
\label{fig:RL2mes}
\end{figure}

Auf der linken Seite wird der Testpin 1 und auf der rechten Seite der Testpin 3 gezeigt.
In beiden Schaltungen kann man erkennen, dass der Anschluss 3 (TP3) mit VCC und die linke Seite (TP1) mit
GND verbunden ist.
Die Stromrichtung durch den Widerstand Rx ist immer die gleiche.
Die Werte f\"ur auf Ausgang geschaltete Ports werden mit roter Farbe dargestellt, 
die Werte f\"ur die Eing\"ange werden mit blauer Farbe dargestellt, inaktive Ports sind schwarz.
In beiden gezeigten Messmethoden sollte der Strom den gleichen Wert haben, weil die Summe der Widerst\"ande zwischen
VCC und GND gleich ist, vorausgesetzt die eingebauten Widerst\"ande sind gleich.
Normalerweise sind die gemessenen Spannungen aber nicht gleich, weil die Reihenfolge
der Widerst\"ande vertauscht ist.

Das V-Symbol innerhalb eines Kreises markiert die Ports, die f\"ur die Spannungsmessung benutzt werden.
In beiden Konfigurationen kann der Wert des Widerstandes Rx aus den bekannten Widerstandswerten
und den gemessenen Spannungen berechnet werden, wenn das Verh\"altnis des Widerstands Rx und den \(680\Omega\)-Widerst\"anden
 nicht zu hoch ist.
Der theoretische Spannungsverlauf wird in Abbildung~\ref{fig:RLvtot} gezeigt, wobei die Widerstandswerte 
in logarithmischer Skalierung dargestellt sind.
\begin{figure}[H]
\centering
\input{../GNU/RLvtot}
\caption{Spannung von Type 1 und Type 2 mit Messwiderstand \(680\Omega\) }
\label{fig:RLvtot}
\end{figure}
Die Verlauf f\"ur Messung Type 1 wird in Abbildung~\ref{fig:RLvlow} mit gespreizter Darstellung f\"ur die unteren Widerstandwerte gezeigt.
Wie man sehen kann, braucht man eine bessere ADC-Aufl\"osung als die möglichen 4,9mV bei der 5V ADC-Referenz, um richtige
Widerstandswerte von der gemessenen Spannung unter \(2\Omega\) zu erhalten.
Es gibt nur drei ADC-Stufen mit der 5V-Referenz zwischen \(0\Omega\) und \(2\Omega\).
Hier kann die Bereichsumschaltung mit der AUTOSCALE\_ADC Option helfen.
Der gleiche gespreizte Bereich f\"ur die Type 2 Messung wird in Abbildung~\ref{fig:RLvhigh} gezeigt.
Ungl\"ucklicherweise kann man nicht die h\"ohere ADC-Aufl\"osung f\"ur Messmethode Type 2 benutzen,
weil die Spannung zu hoch ist und unser ATmega keine differenziellen ADC-Eing\"ange besitzt.
Die Messungen mit den \(680\Omega\)-Widerst\"anden werden bis zu einem Widerstandswert von 
\(20k\Omega\) (Spannung ist unter 169mV) zur Bestimmung des Messergebnisses verwendet.

F\"ur h\"ohere Widerstandswerte werden Messungen mit den \(470k\Omega\)-Widerst\"anden benutzt.
Der Mittelwert von beiden Messungen wird f\"ur den angezeigten Widerstandwert benutzt, wenn alle Messungen ergeben,
dass es kein anderes Bauteil ist.
Wenn die AUTOSCALE\_ADC Funktion benutzt wird und eine der gemessenen Spannungen f\"ur beide Versionen unter 0,98V liegt,
wird ein gewichteter Mittelwert mit Faktor vier f\"ur die Messung mit der Spannung unter 0,98V benutzt. Der andere Wert wird mit Faktor eins bewertet.
Das wird wegen der Faktor vier besseren Aufl\"osung dieser Messung gemacht.
Faktor vier wird nur f\"ur ATmega88, ATmega168 und ATmega328 Prozessoren verwendet, f\"ur ATmega8 wird ein
Faktor zwei als Wichtung benutzt wenn die Spannung unter 0,98V ist, weil die ADC Referenzspannung hier 2,56V statt 1,1V betr"agt.
Wenn der ATmega mehr als 8 KByte Flashspeicher besitzt, wird die Spannungsmessung an den Widerst\"anden so lange verz\"ogert,
bis keine \"Anderung mehr festgestellt wird oder eine Zeitgrenze \"uberschritten wird.
Durch diese Ma"snahme werden auch gro"se Kondensatoren nicht mehr irrt\"umlich als
Widerst\"ande erkannt und der Gleichstrom-Widerstand gro"ser Induktivit\"aten wird richtig gemessen.

\begin{figure}[H]
  \begin{subfigure}[b]{9cm}
    \centering
    \resizebox{9cm}{!}{\input{../GNU/RLvlow}}
    \caption{Type 1 Messung}
    \label{fig:RLvlow}
  \end{subfigure}
  ~
  \begin{subfigure}[b]{9cm}
    \centering
    \resizebox{9cm}{!}{% GNUPLOT: LaTeX picture with Postscript
\begingroup
  \makeatletter
  \providecommand\color[2][]{%
    \GenericError{(gnuplot) \space\space\space\@spaces}{%
      Package color not loaded in conjunction with
      terminal option `colourtext'%
    }{See the gnuplot documentation for explanation.%
    }{Either use 'blacktext' in gnuplot or load the package
      color.sty in LaTeX.}%
    \renewcommand\color[2][]{}%
  }%
  \providecommand\includegraphics[2][]{%
    \GenericError{(gnuplot) \space\space\space\@spaces}{%
      Package graphicx or graphics not loaded%
    }{See the gnuplot documentation for explanation.%
    }{The gnuplot epslatex terminal needs graphicx.sty or graphics.sty.}%
    \renewcommand\includegraphics[2][]{}%
  }%
  \providecommand\rotatebox[2]{#2}%
  \@ifundefined{ifGPcolor}{%
    \newif\ifGPcolor
    \GPcolortrue
  }{}%
  \@ifundefined{ifGPblacktext}{%
    \newif\ifGPblacktext
    \GPblacktexttrue
  }{}%
  % define a \g@addto@macro without @ in the name:
  \let\gplgaddtomacro\g@addto@macro
  % define empty templates for all commands taking text:
  \gdef\gplbacktext{}%
  \gdef\gplfronttext{}%
  \makeatother
  \ifGPblacktext
    % no textcolor at all
    \def\colorrgb#1{}%
    \def\colorgray#1{}%
  \else
    % gray or color?
    \ifGPcolor
      \def\colorrgb#1{\color[rgb]{#1}}%
      \def\colorgray#1{\color[gray]{#1}}%
      \expandafter\def\csname LTw\endcsname{\color{white}}%
      \expandafter\def\csname LTb\endcsname{\color{black}}%
      \expandafter\def\csname LTa\endcsname{\color{black}}%
      \expandafter\def\csname LT0\endcsname{\color[rgb]{1,0,0}}%
      \expandafter\def\csname LT1\endcsname{\color[rgb]{0,1,0}}%
      \expandafter\def\csname LT2\endcsname{\color[rgb]{0,0,1}}%
      \expandafter\def\csname LT3\endcsname{\color[rgb]{1,0,1}}%
      \expandafter\def\csname LT4\endcsname{\color[rgb]{0,1,1}}%
      \expandafter\def\csname LT5\endcsname{\color[rgb]{1,1,0}}%
      \expandafter\def\csname LT6\endcsname{\color[rgb]{0,0,0}}%
      \expandafter\def\csname LT7\endcsname{\color[rgb]{1,0.3,0}}%
      \expandafter\def\csname LT8\endcsname{\color[rgb]{0.5,0.5,0.5}}%
    \else
      % gray
      \def\colorrgb#1{\color{black}}%
      \def\colorgray#1{\color[gray]{#1}}%
      \expandafter\def\csname LTw\endcsname{\color{white}}%
      \expandafter\def\csname LTb\endcsname{\color{black}}%
      \expandafter\def\csname LTa\endcsname{\color{black}}%
      \expandafter\def\csname LT0\endcsname{\color{black}}%
      \expandafter\def\csname LT1\endcsname{\color{black}}%
      \expandafter\def\csname LT2\endcsname{\color{black}}%
      \expandafter\def\csname LT3\endcsname{\color{black}}%
      \expandafter\def\csname LT4\endcsname{\color{black}}%
      \expandafter\def\csname LT5\endcsname{\color{black}}%
      \expandafter\def\csname LT6\endcsname{\color{black}}%
      \expandafter\def\csname LT7\endcsname{\color{black}}%
      \expandafter\def\csname LT8\endcsname{\color{black}}%
    \fi
  \fi
  \setlength{\unitlength}{0.0500bp}%
  \begin{picture}(7200.00,5040.00)%
    \gplgaddtomacro\gplbacktext{%
      \csname LTb\endcsname%
      \put(1078,704){\makebox(0,0)[r]{\strut{} 4780}}%
      \csname LTb\endcsname%
      \put(1078,995){\makebox(0,0)[r]{\strut{} 4785}}%
      \csname LTb\endcsname%
      \put(1078,1286){\makebox(0,0)[r]{\strut{} 4790}}%
      \csname LTb\endcsname%
      \put(1078,1576){\makebox(0,0)[r]{\strut{} 4795}}%
      \csname LTb\endcsname%
      \put(1078,1867){\makebox(0,0)[r]{\strut{} 4800}}%
      \csname LTb\endcsname%
      \put(1078,2158){\makebox(0,0)[r]{\strut{} 4805}}%
      \csname LTb\endcsname%
      \put(1078,2449){\makebox(0,0)[r]{\strut{} 4810}}%
      \csname LTb\endcsname%
      \put(1078,2740){\makebox(0,0)[r]{\strut{} 4815}}%
      \csname LTb\endcsname%
      \put(1078,3030){\makebox(0,0)[r]{\strut{} 4820}}%
      \csname LTb\endcsname%
      \put(1078,3321){\makebox(0,0)[r]{\strut{} 4825}}%
      \csname LTb\endcsname%
      \put(1078,3612){\makebox(0,0)[r]{\strut{} 4830}}%
      \csname LTb\endcsname%
      \put(1078,3903){\makebox(0,0)[r]{\strut{} 4835}}%
      \csname LTb\endcsname%
      \put(1078,4193){\makebox(0,0)[r]{\strut{} 4840}}%
      \csname LTb\endcsname%
      \put(1078,4484){\makebox(0,0)[r]{\strut{} 4845}}%
      \csname LTb\endcsname%
      \put(1078,4775){\makebox(0,0)[r]{\strut{} 4850}}%
      \csname LTb\endcsname%
      \put(1210,484){\makebox(0,0){\strut{} 0}}%
      \csname LTb\endcsname%
      \put(1769,484){\makebox(0,0){\strut{} 1}}%
      \csname LTb\endcsname%
      \put(2329,484){\makebox(0,0){\strut{} 2}}%
      \csname LTb\endcsname%
      \put(2888,484){\makebox(0,0){\strut{} 3}}%
      \csname LTb\endcsname%
      \put(3447,484){\makebox(0,0){\strut{} 4}}%
      \csname LTb\endcsname%
      \put(4007,484){\makebox(0,0){\strut{} 5}}%
      \csname LTb\endcsname%
      \put(4566,484){\makebox(0,0){\strut{} 6}}%
      \csname LTb\endcsname%
      \put(5125,484){\makebox(0,0){\strut{} 7}}%
      \csname LTb\endcsname%
      \put(5684,484){\makebox(0,0){\strut{} 8}}%
      \csname LTb\endcsname%
      \put(6244,484){\makebox(0,0){\strut{} 9}}%
      \csname LTb\endcsname%
      \put(6803,484){\makebox(0,0){\strut{} 10}}%
      \put(176,2739){\rotatebox{-270}{\makebox(0,0){\strut{}voltage / mV}}}%
      \put(4006,154){\makebox(0,0){\strut{}resistor Rx / Ohm}}%
      \put(4006,4665){\makebox(0,0){\strut{}}}%
      \put(286,110){\makebox(0,0)[l]{\strut{}}}%
    }%
    \gplgaddtomacro\gplfronttext{%
      \csname LTb\endcsname%
      \put(5816,4602){\makebox(0,0)[r]{\strut{}PC0, type 2}}%
    }%
    \gplbacktext
    \put(0,0){\includegraphics{../GNU/RLvhigh}}%
    \gplfronttext
  \end{picture}%
\endgroup
}
    \caption{Type 2 Messung}
    \label{fig:RLvhigh}
  \end{subfigure}
  \caption{Ausschnitt des theoretischen Spannungsverlauf von \(0\Omega\) bis \(10\Omega\)}
\end{figure}


\subsection{Widerstandsmessung mit den 470-kOhm-Widerst\"anden}
Die n\"achsten Abbildungen~\ref{fig:RH1mes} und \ref{fig:RH2mes} zeigen die gleichen Messmethoden f\"ur die Messungen mit
 \(470k\Omega\)-Pr\"azisionswiderst\"anden.
Weil \(470k\Omega\) in Relation zu den Port-Widerst\"anden \(22\Omega\) und \(19\Omega\) sehr gro"s ist ,
k\"onnen die Port-Widerst\"ande f\"ur die Berechnung des Widerstandswertes Rx vernachl\"assigt werden.

F\"ur beide Messmethoden mit den \(470k\Omega\)-Widerst\"anden wird nur eine Spannung gemessen, weil der Strom
so niedrig ist, dass keine Spannungsdifferenz an den internen Port Widerst\"anden gemessen werden kann (wie zu erwarten).
Der theoretische Spannungsverlauf wird in Abbildung~\ref{fig:RHv} gezeigt, wobei die Widerstandswerte wieder in
logarithmischer Skalierung gezeigt werden.
Der theoretische Verlauf in diesem Diagramm endet bei \(100M\Omega\), aber das Ergebnis des Testers wird auf
 \(60M\Omega\) begrenzt, anderenfalls nimmt der Tester an, dass kein Widerstand angeschlossen ist.
Als Ergebnis wird der gewichtete Mittelwert von beiden Messmethoden verwendet. Dies geschieht nach den gleichen Regeln, die schon bei
den Messungen mit den  \(680\Omega\) Widerst\"anden beschrieben wurden.
Ich habe beobachtet, dass die Messergebnisse f\"ur alle ATmega Typen  n\"aher am wahren Wert liegen, wenn zum Messergebnis
ein konstanter Offset von \(700\Omega\) addiert wird.
Dieser Offset kann mit der Konstante RH\_OFFSET (define) in der Datei config.h angepasst werden.

\begin{figure}[H]
\centering
\includegraphics[]{../FIG/ResistormessH1.eps}
\caption{Messung Type 3 mit \(470k\Omega\) }
\label{fig:RH1mes}
\end{figure}

\begin{figure}[H]
 \centering
 \includegraphics[]{../FIG/ResistormessH2.eps}
 \caption{Messung Type 4 mit \(470k\Omega\) }
\label{fig:RH2mes}
\end{figure}

\begin{figure}[H]
\centering
\input{../GNU/RHv}
\caption{Spannungen bei Messungen von Type 3 und Type 4  mit \(470k\Omega\) }
\label{fig:RHv}
\end{figure}

\subsection{Ergebnisse der Widerstands-Messung}
Abbildung~\ref{fig:mega8res} zeigt den relativen Fehler der Widerstandmessung mit einem ATmega8. 
,,Mega8'' sind die Ergebnisse ohne AUTOSCALE\_ADC Option,
,,Mega8as'' sind die Ergebnisse mit der Option AUTOSCALE\_ADC und
,,Mega8orig'' sind die gleichen Messungen mit der Originalsoftware von Markus F. ohne meine Verbesserungen.
Die Ergebnisse f\"ur die AUTOSCALE\_ADC Version zeigen hier kaum Verbesserungs-Potential.
Abbildung~\ref{fig:mega168res} zeigt die gleichen Messungen mit einem ATmega168.
,,Mega168'' sind die Ergebnisse ohne die AUTOSCALE\_ADC Option, ,,Mega168as'' die mit der
 AUTOSCALE\_ADC Option.

Mit dem ATmega168 scheint es m\"oglich zu sein, Messungen von Widerst\"anden im
Bereich von \(20\Omega\) bis \(20M\Omega\) mit einem Messfehler von unter \(\pm1\%\) durchzuf\"uhren.
F\"ur Messungen unterhalb von \(100\Omega\) sollte man ber\"ucksichtigen, dass jede Pr\"ufklemme mit Kabel ebenfalls
einen Widerstandwert hat.
Es ist besser, den Widerstand direkt mit den Anschlussklemmen zu verbinden.
Wenn das nicht m\"oglich ist, sollte man der Widerstandswert der kurzgeschlossenen Pr\"ufklemmen vom Messergebnis abziehen.
Zeigt beispielsweise der Tester einen Wert von \(30,6\Omega\) an, wenn der Präzisionswiderstand einen aufgedruckten Wert von \(30\Omega\) hat, 
und die im Kurzschluss gemessenen Pr\"ufklemmen haben einen Wert von \(0,5\Omega\), dann wird der Widerstand vom
Tester mit \(30,1\Omega\) gemessen.
Unterhalb von einem Widerstandswert von \(10\Omega\) macht ein Aufl\"osungs-Schritt von \(0,1\Omega\) schon einen Fehler von mehr als 1\%!

\begin{figure}[H]
\centering
\input{../GNU/Mega8res}
\caption{Relativer Fehler in \% f\"ur Widerstands-Messungen mit ATmega8 }
\label{fig:mega8res}
\end{figure}

\begin{figure}[H]
\centering
\input{../GNU/Mega168res}
\caption{Relativer Fehler in \% f\"ur Widerstand-Messungen mit ATmega168 }
\label{fig:mega168res}
\end{figure}

Das Diagramm \ref{fig:m168res_all} zeigt die Messfehler von drei ATmega168-Prozessoren vor der Kalibration als Punkte, nach der
Kalibration als Linie. Entsprechend werden die Messfehler von drei ATmega168A in Abbildung \ref{fig:m168ares_all} und
die Messfehler von drei ATmega328P in Abbildung \ref{fig:m168pres_all} gezeigt.
Die Messfehler der ATmega328 werden in den Abbildungen \ref{fig:m328res_all} und \ref{fig:m328pres_all} gezeigt.
Nach der automatischen Kalibration bleibt der Messfehler mit einer Ausnahme (ATmega328P-13, \(22 k\Omega\)) im Widerstandsbereich
\(10~\Omega~-~20 M\Omega\) im Bereich \(\pm1~\%\).
Vor der Kalibration k\"onnen die Messfehler bei einigen Prozessoren bis \(\pm~3\%\) betragen.
Der Fehler wird verursacht durch die AUTOSCALE\_ADC-Umschaltung der ADC-Referenz.
Durch den direkten Vergleich einer Kondensatorspannung von unter 1~V, einmal mit der VCC-Referenz und nochmal mit
der internen ADC-Referenz gemessen, kann der Fehler ausgeglichen werden.
In diesem Fall wird die Spannung mit dem gleichen Multiplexer-Kanal gemessen und die Bandgap-Referenz ist auf den AREF-Pin
aufgeschaltet.
Die direkte Vermessung der Bandgap-Referenz durch die direkte Wahl des Multiplexer-Eingangs f\"uhrt leider zu diesem Offset,
der entweder manuell mit der Option REF\_R\_KORR oder automatisch mit der Option AUTO\_CAL des Selbsttestes beseitigt werden kann.
Im AUTO\_CAL-Modus ist REF\_R\_KORR ein zus\"atzlicher Offset zur automatisch gefundenen Spannungsdifferenz.

\begin{figure}[H]
  \begin{subfigure}[b]{9cm}
    \centering
    \resizebox{9cm}{!}{% GNUPLOT: LaTeX picture with Postscript
\begingroup
  \makeatletter
  \providecommand\color[2][]{%
    \GenericError{(gnuplot) \space\space\space\@spaces}{%
      Package color not loaded in conjunction with
      terminal option `colourtext'%
    }{See the gnuplot documentation for explanation.%
    }{Either use 'blacktext' in gnuplot or load the package
      color.sty in LaTeX.}%
    \renewcommand\color[2][]{}%
  }%
  \providecommand\includegraphics[2][]{%
    \GenericError{(gnuplot) \space\space\space\@spaces}{%
      Package graphicx or graphics not loaded%
    }{See the gnuplot documentation for explanation.%
    }{The gnuplot epslatex terminal needs graphicx.sty or graphics.sty.}%
    \renewcommand\includegraphics[2][]{}%
  }%
  \providecommand\rotatebox[2]{#2}%
  \@ifundefined{ifGPcolor}{%
    \newif\ifGPcolor
    \GPcolortrue
  }{}%
  \@ifundefined{ifGPblacktext}{%
    \newif\ifGPblacktext
    \GPblacktexttrue
  }{}%
  % define a \g@addto@macro without @ in the name:
  \let\gplgaddtomacro\g@addto@macro
  % define empty templates for all commands taking text:
  \gdef\gplbacktext{}%
  \gdef\gplfronttext{}%
  \makeatother
  \ifGPblacktext
    % no textcolor at all
    \def\colorrgb#1{}%
    \def\colorgray#1{}%
  \else
    % gray or color?
    \ifGPcolor
      \def\colorrgb#1{\color[rgb]{#1}}%
      \def\colorgray#1{\color[gray]{#1}}%
      \expandafter\def\csname LTw\endcsname{\color{white}}%
      \expandafter\def\csname LTb\endcsname{\color{black}}%
      \expandafter\def\csname LTa\endcsname{\color{black}}%
      \expandafter\def\csname LT0\endcsname{\color[rgb]{1,0,0}}%
      \expandafter\def\csname LT1\endcsname{\color[rgb]{0,1,0}}%
      \expandafter\def\csname LT2\endcsname{\color[rgb]{0,0,1}}%
      \expandafter\def\csname LT3\endcsname{\color[rgb]{1,0,1}}%
      \expandafter\def\csname LT4\endcsname{\color[rgb]{0,1,1}}%
      \expandafter\def\csname LT5\endcsname{\color[rgb]{1,1,0}}%
      \expandafter\def\csname LT6\endcsname{\color[rgb]{0,0,0}}%
      \expandafter\def\csname LT7\endcsname{\color[rgb]{1,0.3,0}}%
      \expandafter\def\csname LT8\endcsname{\color[rgb]{0.5,0.5,0.5}}%
    \else
      % gray
      \def\colorrgb#1{\color{black}}%
      \def\colorgray#1{\color[gray]{#1}}%
      \expandafter\def\csname LTw\endcsname{\color{white}}%
      \expandafter\def\csname LTb\endcsname{\color{black}}%
      \expandafter\def\csname LTa\endcsname{\color{black}}%
      \expandafter\def\csname LT0\endcsname{\color{black}}%
      \expandafter\def\csname LT1\endcsname{\color{black}}%
      \expandafter\def\csname LT2\endcsname{\color{black}}%
      \expandafter\def\csname LT3\endcsname{\color{black}}%
      \expandafter\def\csname LT4\endcsname{\color{black}}%
      \expandafter\def\csname LT5\endcsname{\color{black}}%
      \expandafter\def\csname LT6\endcsname{\color{black}}%
      \expandafter\def\csname LT7\endcsname{\color{black}}%
      \expandafter\def\csname LT8\endcsname{\color{black}}%
    \fi
  \fi
  \setlength{\unitlength}{0.0500bp}%
  \begin{picture}(7200.00,5040.00)%
    \gplgaddtomacro\gplbacktext{%
      \csname LTb\endcsname%
      \put(682,704){\makebox(0,0)[r]{\strut{}-5}}%
      \csname LTb\endcsname%
      \put(682,1111){\makebox(0,0)[r]{\strut{}-4}}%
      \csname LTb\endcsname%
      \put(682,1518){\makebox(0,0)[r]{\strut{}-3}}%
      \csname LTb\endcsname%
      \put(682,1925){\makebox(0,0)[r]{\strut{}-2}}%
      \csname LTb\endcsname%
      \put(682,2332){\makebox(0,0)[r]{\strut{}-1}}%
      \csname LTb\endcsname%
      \put(682,2740){\makebox(0,0)[r]{\strut{} 0}}%
      \csname LTb\endcsname%
      \put(682,3147){\makebox(0,0)[r]{\strut{} 1}}%
      \csname LTb\endcsname%
      \put(682,3554){\makebox(0,0)[r]{\strut{} 2}}%
      \csname LTb\endcsname%
      \put(682,3961){\makebox(0,0)[r]{\strut{} 3}}%
      \csname LTb\endcsname%
      \put(682,4368){\makebox(0,0)[r]{\strut{} 4}}%
      \csname LTb\endcsname%
      \put(682,4775){\makebox(0,0)[r]{\strut{} 5}}%
      \csname LTb\endcsname%
      \put(814,484){\makebox(0,0){\strut{}1 }}%
      \csname LTb\endcsname%
      \put(1563,484){\makebox(0,0){\strut{}10 }}%
      \csname LTb\endcsname%
      \put(2311,484){\makebox(0,0){\strut{}100 }}%
      \csname LTb\endcsname%
      \put(3060,484){\makebox(0,0){\strut{}1k}}%
      \csname LTb\endcsname%
      \put(3809,484){\makebox(0,0){\strut{}10k}}%
      \csname LTb\endcsname%
      \put(4557,484){\makebox(0,0){\strut{}100k}}%
      \csname LTb\endcsname%
      \put(5306,484){\makebox(0,0){\strut{}1M}}%
      \csname LTb\endcsname%
      \put(6054,484){\makebox(0,0){\strut{}10M}}%
      \csname LTb\endcsname%
      \put(6803,484){\makebox(0,0){\strut{}100M}}%
      \put(176,2739){\rotatebox{-270}{\makebox(0,0){\strut{}Error / Percent}}}%
      \put(3808,154){\makebox(0,0){\strut{}Resistor value / Ohm}}%
    }%
    \gplgaddtomacro\gplfronttext{%
      \csname LTb\endcsname%
      \put(5753,4602){\makebox(0,0)[r]{\strut{}m168-1}}%
      \csname LTb\endcsname%
      \put(5753,4382){\makebox(0,0)[r]{\strut{}m168-2}}%
      \csname LTb\endcsname%
      \put(5753,4162){\makebox(0,0)[r]{\strut{}m168-3}}%
      \csname LTb\endcsname%
      \put(5753,3942){\makebox(0,0)[r]{\strut{}m168-1}}%
      \csname LTb\endcsname%
      \put(5753,3722){\makebox(0,0)[r]{\strut{}m168-2}}%
      \csname LTb\endcsname%
      \put(5753,3502){\makebox(0,0)[r]{\strut{}m168-3}}%
    }%
    \gplbacktext
    \put(0,0){\includegraphics{../GNU/m168res_all}}%
    \gplfronttext
  \end{picture}%
\endgroup
}
    \caption{mit drei ATmega168}
    \label{fig:m168res_all}
  \end{subfigure}
  ~
  \begin{subfigure}[b]{9cm}
    \centering
    \resizebox{9cm}{!}{\input{../GNU/m168ares_all}}
    \caption{mit drei ATmega168A}
    \label{fig:m168ares_all}
  \end{subfigure}
\caption{Relativer Fehler in \% f\"ur Widerstands-Messungen}
\end{figure}

\begin{figure}[H]
\centering
\input{../GNU/m168pres_all}
\caption{Relativer Fehler in \% f\"ur Widerstands-Messungen mit drei ATmega168P }
\label{fig:m168pres_all}
\end{figure}

\begin{figure}[H]
  \begin{subfigure}[b]{9cm}
    \centering
    \resizebox{9cm}{!}{\input{../GNU/m328res_all}}
    \caption{mit drei ATmega328}
    \label{fig:m328res_all}
  \end{subfigure}
  ~
  \begin{subfigure}[b]{9cm}
    \centering
    \resizebox{9cm}{!}{\input{../GNU/m328pres_all}}
    \caption{mit drei ATmega328P}
    \label{fig:m328pres_all}
  \end{subfigure}
\caption{Relativer Fehler in \% f\"ur Widerstands-Messungen}
\end{figure}

